\chapter{Literature Review}

\section{Introduction}

Microservice architectures have become the dominant paradigm for building scalable, cloud-native applications. However, the distributed nature of microservices introduces significant challenges in fault detection and root cause analysis (RCA). When a fault occurs in a multi-service system, identifying which component is responsible requires sophisticated analysis of telemetry data including metrics, logs, and distributed traces. This literature review examines the current state of research in microservice fault detection and localisation, with particular focus on the trade-offs between deep learning approaches and lightweight rule-based methods.

\section{Deep Learning Approaches to Microservice Anomaly Detection}

Recent advances in deep learning have produced highly accurate microservice fault detection systems, but at the cost of substantial computational resources and training data requirements.

\subsection{Multi-Modal Deep Learning Methods}

\citet{xing2025dualchannel} present a state-of-the-art dual-channel deep learning approach combining Temporal Convolutional Networks (TCN) and Variational Autoencoders (VAE) with contrastive learning and causal inference for multi-dimensional anomaly detection in microservice architectures. Their method achieves an impressive detection F1-score of 93.8\%, accuracy of 95.4\%, and fault-component localisation precision of 87.6\% under semi-supervised evaluation conditions. This work represents the current accuracy ceiling for learned methods when sufficient labelled training data is available. However, their approach requires extensive model training on large telemetry corpora, making it inapplicable to newly deployed serverless services with no historical data.

\citet{zhang2026cpuonly} propose a CPU-only spatiotemporal anomaly detection system using dynamic graph neural networks (GNN) and LSTM for microservice environments. Their work demonstrates that effective deep learning detection can be achieved without GPU acceleration, though training still requires substantial computational resources and historical telemetry. The temporal modeling via LSTM captures service call patterns over time, while the GNN component encodes service dependencies. This architecture is particularly relevant for understanding how graph-based representations can improve localisation accuracy beyond simple metric thresholds.

\citet{zhang2026galr} introduce GALR (Graph-based root cause Localization and LLM-Assisted Recovery), which combines graph-based root cause localisation with large language model assistance for automated recovery in microservice systems. Their graph construction from service dependencies enables more accurate fault propagation analysis than flat metric-based approaches. However, the LLM-assisted recovery component adds deployment complexity and inference latency that may be prohibitive in cost-constrained serverless environments.

\subsection{End-to-End Troubleshooting Frameworks}

\citet{lee2023eadro} present Eadro, an end-to-end troubleshooting framework that jointly performs detection and localisation on multi-source data (metrics, logs, traces). Eadro represents an important step toward practical deployment by avoiding separate training of detection and localisation components. However, their evaluation shows that the joint approach still requires substantial training data from normal and anomalous periods, limiting applicability to greenfield deployments.

\citet{zhang2022deeptralog} combine trace and log data for microservice anomaly detection through graph-based deep learning in their DeepTraLog system. By constructing execution graphs from traces and enriching them with log embeddings, they achieve strong detection performance. Their work highlights the value of multi-modal telemetry fusion, but the graph construction and embedding training add preprocessing overhead that compounds the training data requirements.

\subsection{Trace-Based Root Cause Localisation}

\citet{zhang2024tracecontrast} present TraceContrast, a trace-based multi-dimensional root cause localisation method for performance issues in microservice systems. Their contrastive learning approach learns to distinguish normal and anomalous trace patterns, achieving effective localisation without requiring explicit fault labels. However, the method still requires a corpus of traces from normal operation to learn the contrastive representations, and the inference time grows with the number of services.

\citet{yu2023nezha} introduce Nezha, an interpretable fine-grained root cause analysis system using multi-modal observability data. Nezha's key contribution is providing human-readable explanations for localisation decisions by attributing anomalies to specific metrics and service call patterns. This interpretability is valuable for operations teams, but comes at the cost of additional feature engineering and model complexity.

\section{Causal Inference and Unsupervised Root Cause Analysis}

A parallel line of research applies causal inference and unsupervised learning to avoid the labeled training data requirements of supervised deep learning.

\subsection{Bayesian and Causal Methods}

\citet{pham2024baro} present BARO (Bayesian Online Change Point Detection), a robust root cause analysis method that uses multivariate Bayesian online change point detection to identify when service metrics deviate from normal patterns. BARO's unsupervised approach avoids training data requirements, but relies on accurate statistical modeling assumptions that may not hold in highly variable serverless workloads. Their evaluation on the RCAEval benchmark demonstrates competitive performance against supervised methods.

\citet{li2022circa} introduce CIRCA (Causal Inference-based Root Cause Analysis), which uses intervention recognition and causal graphs to localise faults in online service systems. CIRCA builds a causal graph from service dependencies and uses perturbation analysis to identify root causes. While CIRCA avoids supervised training, it requires substantial computational resources to construct and analyse the causal graph for each incident, with reported analysis times of several minutes per case.

\citet{xin2023causalrca} propose CausalRCA, a causal inference-based approach for precise fine-grained root cause localisation in microservices. CausalRCA uses graph neural networks with PageRank-based scoring to identify root causes from service dependency graphs. Published precision can be strong (literature Top-1 figures are citation-only here), but GNN inference adds latency that may be prohibitive for real-time alerting in serverless environments. In this artefact, CausalRCA is incomplete ($n{=}4$) and fixed-order flagged, so it is quarantined from peer ranking.

\subsection{TraceRCA-Style Spectrum-Based Localisation}

\citet{li2021tracerca} present the foundational TraceRCA method for practical root cause localisation via trace analysis. TraceRCA uses spectrum-based fault localisation adapted from software debugging: services are ranked by how frequently they appear in failed traces versus successful traces. This lightweight approach requires no training and runs in milliseconds, making it suitable for real-time operation. However, TraceRCA's accuracy suffers when multiple services exhibit correlated anomalies or when fault propagation patterns are complex. TraceRCA serves as an important baseline for understanding the accuracy floor of training-free methods.

\section{Training-Less and Statistical Anomaly Detection}

Recognising the deployment barriers of deep learning methods, several recent works focus on training-less or lightweight statistical approaches.

\subsection{Threshold-Based and Statistical Methods}

\citet{elkhairi2024replicawatcher} introduce ReplicaWatcher, a training-less anomaly detection system for containerised microservices. ReplicaWatcher monitors container replica behaviour and flags anomalies when metric distributions deviate significantly from the ensemble. Their approach demonstrates that effective detection is possible without historical training data by leveraging spatial redundancy across replicas. However, ReplicaWatcher focuses on detection and does not address root cause localisation, which is critical for remediation.

\citet{liu2023practical} present practical methods for anomaly detection over multivariate monitoring metrics using adaptive thresholding. Their work validates that simple statistical approaches (e.g., 3-sigma rules, quantile-based thresholds) remain competitive with complex models when tuned carefully. However, they note that fixed thresholds suffer from concept drift and recommend periodic recalibration—a challenge in serverless environments where traffic patterns change rapidly.

\citet{schmidl2022anomaly} provide a comprehensive evaluation of time series anomaly detection methods, comparing classical statistical approaches against deep learning variants. Their findings show that no single method dominates across all datasets and anomaly types. For high-noise, limited-data scenarios (typical of new serverless deployments), they find that robust statistical methods often outperform deep learning due to reduced overfitting.

\section{Benchmarks and Evaluation Frameworks}

The lack of standardised benchmarks has historically made it difficult to compare microservice RCA methods. Recent efforts address this gap.

\subsection{RCAEval: A Standard Benchmark}

\citet{pham2025rcaeval} introduce RCAEval, a comprehensive benchmark for root cause analysis containing 735 labelled failure cases across three microservice systems (Online Boutique, Sock Shop, Train Ticket). RCAEval provides metrics, logs, and traces for each case, along with implementations of 15 reproducible baseline methods including TraceRCA, BARO, and CIRCA variants. The benchmark defines standard evaluation metrics (AC@k, Avg@k) and provides a harness for fair comparison. RCAEval represents a significant advance in reproducibility for microservice RCA research and enables direct comparison of the rule-based approach proposed in this work against published baselines.

\subsection{Domain-Specific Datasets}

\citet{hardt2024petshop} present the PetShop dataset for root cause analysis of performance issues in microservices. PetShop focuses specifically on performance degradation rather than binary fault/no-fault scenarios, providing fine-grained labels for latency anomalies. Their work demonstrates that different fault types (crashes, latency spikes, resource exhaustion) may require different localisation strategies, motivating the multi-fault-type evaluation in this research.

\section{Monitoring Overhead and Sampling Strategies}

An often-overlooked dimension in microservice RCA research is the cost and overhead of telemetry collection itself.

\subsection{Distributed Tracing Overhead}

\citet{huang2024trastrainer} introduce TraStrainer, an adaptive sampling system for distributed traces that adjusts sampling rates based on system runtime state. They demonstrate that full tracing (100\% sampling) can add 5-15\% end-to-end latency overhead and generate overwhelming telemetry volumes. TraStrainer reduces overhead by sampling more aggressively during normal operation and increasing coverage during suspected anomalies. This dynamic approach is promising but adds control-plane complexity.

\citet{mertz2023software} study software runtime monitoring with adaptive sampling rates, showing that representative trace coverage can be achieved with as low as 1-5\% sampling under steady workloads. However, they note that adaptive sampling introduces a cold-start problem: anomalies must be detected before sampling rates can be increased, potentially missing early fault symptoms. This finding motivates the fixed-policy sampling (5\%) used in the overhead evaluation of this work.

\citet{eder2023comparison} compare distributed tracing tools in serverless applications, specifically examining AWS X-Ray, Jaeger, and Zipkin. They find that X-Ray imposes lower baseline latency overhead (~2-5ms) compared to OSS alternatives but offers less flexibility in custom instrumentation. For AWS Lambda deployments, X-Ray's native integration and automatic context propagation make it the practical choice despite its proprietary nature.

\section{AWS-Specific Observability and Serverless Monitoring}

The official AWS documentation provides the ground truth for CloudWatch and X-Ray capabilities in serverless environments.

\subsection{CloudWatch Metrics and Alarms}

\citet{awscloudwatch} describe CloudWatch's built-in Lambda metrics (Invocations, Errors, Duration, Throttles) and alarm capabilities. CloudWatch alarms support static thresholds and anomaly detection models, but the anomaly detection feature requires at least 2 weeks of training data—unavailable in new deployments. This limitation motivates the 3-sigma rule-based detection approach used in this work.

\subsection{X-Ray Tracing and Service Maps}

\citet{awsxray} document X-Ray's active tracing capabilities for Lambda functions and API Gateway. X-Ray automatically captures service call graphs and timing breakdowns, enabling dependency analysis without custom instrumentation. The service map visualization is generated server-side, but the raw trace segments can be retrieved via API for programmatic root cause ranking. This research leverages X-Ray's segment-level data to implement the dependency ranker described in Chapter~5.

\subsection{Lambda-Specific Considerations}

\citet{awslambda} detail Lambda's integration with X-Ray and CloudWatch. Key considerations for serverless fault detection include: cold-start effects (first invocation of a container is slower), concurrent execution limits (throttling occurs when concurrency is exhausted), and the asynchronous nature of Lambda-to-Lambda calls. The fault taxonomy used in this work aligns with these Lambda-specific failure modes.

\subsection{Pricing and Cost Models}

\citet{awscloudwatchpricing, awsxraypricing} specify the pricing for CloudWatch and X-Ray services. Costs are volume-based: CloudWatch charges per metric-month and per custom metric, while X-Ray charges per trace recorded and per trace retrieved. At production scale (millions of requests), tracing costs can exceed compute costs if sampling is not carefully controlled. The overhead evaluation in Chapter~7 includes detailed cost estimates based on these pricing models.

\section{Identified Research Gap}

This literature review reveals a clear gap in the microservice RCA landscape:

\begin{itemize}
    \item \textbf{Deep learning methods} (Xing et al., Zhang et al., Lee et al., Zhang et al.) achieve state-of-the-art accuracy (F1 $>$ 90\%, localisation precision $>$ 85\%) but require large labeled training corpora and substantial compute resources for training and inference. They are inapplicable to newly deployed serverless services.
    
    \item \textbf{Training-free methods} (ReplicaWatcher, TraceRCA) operate without historical data but sacrifice localisation accuracy, often falling below 70\% Top-3 accuracy on complex multi-service failures.
    
    \item \textbf{Causal and Bayesian methods} (CIRCA, BARO, CausalRCA) avoid supervised training but require minutes of processing time per incident due to graph construction and causal analysis—prohibitive for real-time alerting.
    
    \item \textbf{Monitoring overhead} is rarely measured alongside accuracy. Most RCA research reports detection and localisation metrics but ignores the telemetry volume, latency impact, and dollar cost of the proposed monitoring regime.
\end{itemize}

\textbf{This work addresses the gap by:}

\begin{enumerate}
    \item Proposing a \textit{hybrid lightweight approach} that combines fast rule-based detection (CloudWatch 3-sigma thresholds) with lightweight unsupervised ML (Isolation Forest) for localisation—avoiding both the training data requirements of deep learning and the graph complexity of causal methods.
    
    \item Establishing a \textit{three-leg evaluation protocol} that fairly compares the rule-based approach against (1) the published deep learning ceiling (Xing et al.), (2) reproducible baselines on common benchmark data (RCAEval), and (3) overhead metrics that deep methods cannot provide on the same system.
    
    \item Demonstrating that competitive accuracy (F1 within 10 percentage points of the deep learning ceiling) can be achieved with $\geq$50\% reduction in telemetry overhead—a position on the accuracy-overhead Pareto frontier that prior work has not empirically established.
\end{enumerate}

The next chapter describes the research methodology used to evaluate this hybrid approach across all three dimensions: accuracy on common data, localisation precision, and monitoring overhead in a live AWS serverless deployment (or simulated equivalent where live deployment is infeasible).
